The purpose of the demonstrator is to visualize how tensile and compressive loads effect a drive system when implemented, and what steps are needed, to realize a reliable data transmission throughout the system. As in cooperation with SYDE \cite{SydeTechnology:2025} the second purpose of the demonstrator is to present a model of the \ac{HTA} system \cite{SydeHTA:2021} and to explore how it can be applied in different use cases. It serves as a testing platform to study system’s behavior in practice, when implementing new features.
The focus of this thesis will lay on the data transmission in the \ac{HTA} system, first between the Strain Gauges and the Hypatia Board\cite{SydeHypatia:2023} as control unit, second on how to implement this signal to a drive system. After a here following Hypothesis of where Data can get lost, or in other words, where the accuracy of data transmission leaks, the implementation of the system into a demonstrator is described. When a good overview of all the demonstrators components is given, a discussion is opened, where the strengths and the potential for further optimization of the Demonstrator are laying. 

With this overview in mind, the predefined requirements for the demonstrator and an introduction into the \ac{HTA} system as well as into the general architecture of the demonstrator is provided in the following. 
\newpage